% !TEX root = CSL-main.tex
\section{Progress Domain Condition}
\label{sec:progressDomainCond}

%Given the equivalence between SCT and GTC established in Proposition \ref{PropEqualCond}, it becomes evident that both conditions share two key characteristics: (1) both can be formulated using size-change graphs, and (2) the composition of these graphs relies on an underlying structure. In the case of SCT, this structure arises from the control flow between function calls. For GTC, it is determined by the position of the companion nodes in the cyclic pre-proofs. Motivated by these commonalities, we now introduce a unified abstraction that captures both SCT and GTC within a single formal definition. 

%\subsection{Bi-colored relation transition graphs}
%Building on observation (1), we begin by abstracting size-change graphs as \emph{bi-colored relations} over a finite set.

Let $RR'$ be the ordinary relation composition; we also write $x R y$ for $(x,y) \in R$.

\begin{definition}[Bi-colored relations]
%\textbf{(Bi-colored relations and pair composition)} 
Let $X = \{x_1, ..., x_n\}$ be a finite set of names. A {\em bi-colored relation} over $X$, is a pair $(S,P)$ such that $S, P \subseteq X \times X$ where $S \cap P = \emptyset$. For two names $x,y \in X$, if $x S y$ (resp. $xPy$) we say that $x$ \emph{weakly progresses} (resp. \emph{progresses}) to $y$ in $(S, P)$. 
If $(S,P)$ and $(S',P')$ are bi-colored relations then we define the {\em composition} $(S, P) \circ (S', P') = (S'' \setminus P'',\; P'')$
%\[ (S, P) \circ (S', P') = (S'' \setminus P'',\; P'') \]
where
\[S'' = S S \qquad \mbox{and} \qquad P'' = (S P') \cup (P S') \cup (P P') \]
%We define the composition $(S,P) \circ (S',P')$ of two bi-colored relation $(S,P)$ and $(S',P')$ as:
%    \begin{align*}
%    (S, P) \circ (S', P') &= (S'' \setminus P'',\; P'')
%    \end{align*}
%    where $S''$ and $P''$ are defined as:
%    \begin{align*}
%    S'' = S S  \quad P'' = (S P') \cup (P S') (P \circ P')
%    \end{align*}
\end{definition}

We observe that the composition of bi-colored relations is compatible with the union and ordinary composition of binary relations: let's write 
$|(S,P)| = S \cup P$, then
\begin{align*}
|(S,P) \circ (S',P')| = |(S'' \setminus P'', P'')| &= (S''\setminus P'') \cup P'' = SS' \cup  (S P') \cup (P S') \cup (P P') \\
& = (S \cup P) (S' \cup P') = |(S,P)| |(S',P')|
\end{align*}


\begin{definition}[Transition Graph]
Let $\mathcal{B} = \{(S_1,P_1),...,(S_n,P_n) \}$ be a set of bi-colored relations over the same set $X$ of names. A {\em transition graph} of
$\mathcal{B}$ is a tuple $\mathcal{G} =(V, \mathcal{B}, E)$ where $V$ is a finite set of vertices, and $E \subseteq V \times \mathcal{B}\times V$
is a (finite) set of labeled edges.
\end{definition}


\newpage
%To capture observation (2), we introduce a transition system whose edges are labeled by bi-colored relations, abstracting both the control flow between function calls in SCT and the proof structure in cyclic pre-proofs for GTC. This general notion provides a common foundation for the unified treatment of SCT and GTC.

\begin{definition}[Bi-colored relation transition graph]
% \textbf{(Bi-colored relation transition graph):} 
Let $\mathcal{B} = \{(S_1,P_1),...,(S_n,P_n) \}$ be a set of bi-colored relations over the same set of names. 
We define a \emph{bi-colored relation transition graph}, as a labeled transition system represented by the tuple $\mathcal{T} =(V, \mathcal{B}, E)$ where: $V$ is a finite non-empty set of vertices and $E \subseteq V \times \mathcal{B}\times V$ is a set of directed edges labeled by bi-colored relations in $\mathcal{B}$.   
\end{definition}

\begin{definition}
    \textbf{(Composable)}
\end{definition}

\begin{example}
    \textcolor{red}{[Include here an example (THE SAME AS FOR SECTION 4) of how the three defintions SCT <-> GTC <-> transition graph interact][Make basically a triangle where each corner is each representation and each edge is a biconditional <->]}
\end{example}
